\section{Introduction}\label{sec:nintro}

A topos can be presented in many ways that look unrelated: as the actions of a monoid, as the representations of a quiver, as the sheaves on a space, as the models of a logical theory. O.~Caramello's \emph{bridges} \cite{caramello-unification,caramello-tst} exploit this. A property of the topos that is invariant under equivalence is characterised on each of two presentations; these characterisations are the \emph{arches} of the bridge, and the equivalence between the two presentations is its \emph{deck}. Put together, the two arches give a theorem relating the two presentations, in which the topos no longer appears.

In a talk at Mines Paris \cite{caramello-mines}, Caramello proposed that a machine combine arches and decks systematically. The companion article \cite{bridges-generator} runs a small version of this generator by hand and describes it as a pipeline: a registry of arches, proved in Lean~4 as lemmas of the form \texttt{(property of PSh C) ↔ (condition on C)}; a registry of decks, small categories $C$ whose presheaves are a familiar structure; and stages that pair every deck with every arch, compute the condition on $C$, rewrite it in the language of the field, and check it in Lean, by a direct proof and against the literature. This note applies the generator to Grothendieck's manuscript notes kept at the University of Montpellier, which already contain decks of this kind and invariants read on two presentations of one topos.

\subsection*{Plan}

Section~\ref{sec:recall} recalls the notation and the five arches of the companion article. Section~\ref{sec:morearches} gives three more arches (two-valued, connected, spatial), proved in Lean, and Section~\ref{sec:grdecks} reads them on three decks that recur in Grothendieck's notes. Section~\ref{sec:precedents} collects two precedents in his notes: the exercise of folder 161-3, p.~26, which the generator re-derives and whose preorder form we prove in Lean, with the same identification in folder 37, p.~48; and the subtopoi of a theory in folder 161-2. Section~\ref{sec:pseudofields} adds a deck taken from folder 161-2, the topos of pseudo-fields, which Grothendieck presents in three ways; we check one of the presentations, read the arches on the deck, and answer a question he left in the margin. Section~\ref{sec:stacks} relates the classifying topoi to Grothendieck's stacks and proposes finer arches, with an example from folder 146. Section~\ref{sec:nconclusion} concludes, and Appendix~\ref{sec:nlean} lists the Lean files and declarations.
